%% File: hecke-alg-corresp-for-metapletic-SL2-revised.tex
%% Revised version of "A correspondence for the generalized Hecke algebra of
%% the metaplectic cover \overline{SL(2,F)}, F p-adic", NYJM 4 (1998), 223-235.
%% Draft 1 of the corrected and expanded version (November 2026).
%% The revision log is hecke-alg-corresp-for-metapletic-SL2-revision-log.tex.
%% Compiles cleanly with pdflatex.

\documentclass[12pt,letterpaper,reqno]{amsart}
\usepackage[T1]{fontenc}
\usepackage{newtxtext,newtxmath}
\usepackage[letterpaper,margin=1in]{geometry}
\usepackage{amsmath,amssymb,amsthm}
\usepackage{array,longtable}
\usepackage{microtype}
\usepackage[colorlinks=true,linkcolor=blue,citecolor=blue,urlcolor=blue]{hyperref}

% Legacy AMS-LaTeX compatibility macros
\let\cal\mathcal
\let\Bbb\mathbb
\let\frak\mathfrak

% Theorems, propositions, lemmas, corollaries, definitions, examples and
% remarks share one counter, numbered within sections (style guide 1.6).
\newtheorem{theorem}{Theorem}[section]
\newtheorem{corollary}[theorem]{Corollary}
\newtheorem{conjecture}[theorem]{Conjecture}
\newtheorem{lemma}[theorem]{Lemma}
\newtheorem{proposition}[theorem]{Proposition}

\theoremstyle{definition}
\newtheorem{definition}[theorem]{Definition}
\newtheorem{example}[theorem]{Example}
\newtheorem{axiom}{Axiom}

\theoremstyle{remark}
\newtheorem{remark}[theorem]{Remark}
\newtheorem{exercise}{Exercise}[section]

\numberwithin{equation}{section}

\title[The generalized Hecke algebra of $\overline{SL(2,F)}$]{A correspondence
for the generalized Hecke algebra of the metaplectic cover
$\overline{SL(2,F)}$, $F$ $p$-adic\\[1ex]
{\normalsize Draft 1. Corrected and expanded version}}

\author{David Joyner}
\email{wdjoyner@gmail.com}

\date{Original draft: 1998. Corrected and expanded draft: 2026.}

\thanks{The original version of this paper was published in \emph{New York
J.\ Math.}\ \textbf{4} (1998), 223--235. This corrected and expanded version
(Draft~1) was prepared in November 2026.}

\keywords{metaplectic groups, generalized Hecke algebras,
representation theory, $p$-adic groups, $SL(2)$}
\subjclass[2020]{22E50, 11F70}

\newcommand{\abstracttext}{This corrected and expanded version was prepared
with the help of AI systems (Argon, ChatGPT and Claude). The Acknowledgments
give details. We prove, using a technique developed for
$GL(n)$ in Howe and Moy \cite{HM85}, a bijection between generalized Hecke
algebras of $G=SL(2,F)$ over a $p$-adic field and those of its $n$-fold
metaplectic cover $\overline{G}$. This result implies that there is a
canonical correspondence between irreducible admissible representations of $G$
and genuine irreducible admissible representations of $\overline{G}$ of
``sufficiently large level'' (depending on $n,p$).}

\begin{document}

\begin{abstract}
\abstracttext
\end{abstract}

\maketitle

\tableofcontents

%%%% **** The text of the paper starts here **** %%%%

\section{Introduction}
\label{sec:intro}

This section introduces the covering group $\overline{G}$ studied in the
paper and states what the paper proves. It ends with a running example and a
description of how the paper is organized.

Let $F$ be a $p$-adic field of characteristic $0$, with uniformizer $\pi$,
ring of integers ${\cal O}$, and let $q$ denote the cardinality of the residue
field. Let $\mu_n(F)$ denote the group of $n^{\text{th}}$ roots of unity in $F$
and assume $|\mu_n(F)|=n$. We will identify $\mu_n(F)$ with a subgroup of
${\Bbb{C}}^\times$ (via some fixed isomorphism
$\theta:\mu_n(F)\rightarrow \mu_n({\Bbb{C}})$). Write $\mu_n=\mu_n(F)$.
Let $G=SL(2,F)$ and let $\frak{g}$ denote its Lie algebra. Let
$\overline{G}=\overline{SL(2,F)}$ denote the $n$-fold metaplectic cover
defined by the Kubota cocycle
\[
\beta(g_1,g_2)=\left(\frac{x(g_1g_2)}{x(g_1)},
\frac{x(g_1g_2)}{x(g_2)}\right),
\]
for $g_1,g_2\in G$, where
\[
x\begin{pmatrix}
a & b\\
c & d
\end{pmatrix}=
\begin{cases}
c, & \text{if } c\neq 0,\\
d, & \text{otherwise},
\end{cases}
\]
and where $(\ ,\ )$ is the $n^{\text{th}}$ power Hilbert symbol on
$F^\times \times F^\times$. By \S 1 of Kubota \cite{Ku67}, $\beta$ is a
$2$-cocycle (called a ``factor set'' in \cite{Ku67}) which is trivial near the
identity in $G\times G$, so $\overline{G}$ is a topological covering group of
$G$. We denote elements of $\overline{G}$ by $\overline{x}=(x,\zeta)$,
$x\in G,\ \zeta\in \mu_n$, and multiplication by
\[
\overline{x}_1\overline{x}_2=(x_1,\zeta_1)(x_2,\zeta_2)
=(x_1x_2,\beta(x_1,x_2)\zeta_1\zeta_2).
\]
Let $\rho:\overline{G}\rightarrow G$ denote the natural surjection. If $H$ is
any subgroup of $G$ then we define $\overline{H}=\rho^{-1}(H)$. We choose a
Haar measure on $G$, to be normalized below, and a Haar measure on $\mu_n$
normalized so that $\operatorname{vol}(\mu_n)=1$. Let the Haar measure on
$\overline{G}$ be the product measure.

We prove, using a technique developed for $GL(n)$ in \cite{HM85}, a bijection
between generalized Hecke algebras of $G$ and those of $\overline{G}$. A
precise statement will be given in Theorem \ref{thrm:main}. This result
implies that there is a canonical one-to-one correspondence between
irreducible admissible representations of $G$ and genuine irreducible
admissible representations of $\overline{G}$ of ``sufficiently large level''
(depending on $n,p$). For the precise statement, see Corollary
\ref{corollary:main}.

\begin{example}[The running example]
\label{ex:running}
Throughout the paper we illustrate the definitions with one example. Take
$F=\mathbb{Q}_3$, so that $\pi=3$, ${\cal O}=\mathbb{Z}_3$ and $q=3$. Take
$n=2$, so that $\mu_2=\{\pm 1\}$ and $\overline{G}$ is the classical
two-fold metaplectic cover of $SL(2,\mathbb{Q}_3)$. The quadratic Hilbert
symbol is $(a,b)=1$ when $a$ and $b$ are both units, because $p=3$ is odd. For
a unit $u$ it is $(u,3)=\left(\frac{\bar{u}}{3}\right)$, the Legendre symbol
of the residue $\bar{u}\in\mathbb{F}_3$. For example, $(-1,3)=-1$, since $-1$
is not a square modulo $3$. Here $\gcd(p,2n)=\gcd(3,4)=1$, so every result of
the paper applies, including those of Section \ref{section:savin}. We use
level $k=1$ for the generalized Hecke algebras.
\end{example}

The paper is organized as follows. Section \ref{sec:litreview} reviews the
literature. Section \ref{sec:background} explains the background on Hecke
algebras, covering groups and representations. Section \ref{sec:notation}
fixes notation and defines the two Hecke algebras. Section \ref{sec:lemmas}
proves some lemmas on convolution of genuine functions. Sections
\ref{section:sl2} and \ref{section:metaplectic} give generators and relations
for the two algebras. Section \ref{sec:application} proves the main theorem
and its consequences for representations. Section \ref{section:savin}
discusses the Iwahori case. Section \ref{sec:summary} summarizes the results,
and Section \ref{sec:glossary} is a glossary.

\section{Literature review}
\label{sec:litreview}

This section places the paper among related work. Section \ref{sec:lit-before}
covers work up to the original 1998 version of the paper. Section
\ref{sec:lit-after} covers later work on Hecke algebras of covering groups.

\subsection{Work up to 1998}
\label{sec:lit-before}

The paper builds on four tools. Kubota \cite{Ku67} constructed the $n$-fold
cover of $SL(2,F)$ by an explicit cocycle, which is the cocycle $\beta$ used
here. Kazhdan and Patterson \cite{KP84} studied $n$-fold covers of $GL(r)$
and their genuine representations in general. Howe and Moy \cite{HM85}
determined generators and relations for the generalized Hecke algebras of
$GL(n)$ with respect to filtration subgroups of the Iwahori subgroup. Our
proof follows their method closely. Borel \cite{Bo76} related representations
with Iwahori-fixed vectors to modules over the Iwahori--Hecke algebra, using
the induced and produced modules that we adapt in Section
\ref{sec:application}. Silberger's book \cite{Si79} supplies the facts on
matrix coefficients that we use there.

The study of genuine admissible representations on $p$-adic groups has its
origins in the book \cite{Ge76} by S. Gelbart. Flicker \cite{Fl80} determined a
comparison of Hecke algebras of \emph{smooth functions} on $GL(2,F)$ and
$\overline{GL(2,F)}$ using the matching of orbital integrals. This was then
used, with the Arthur-Selberg trace formula, to establish a global ``Shimura
correspondence''. Neither of these two works used the Hecke algebras
considered here. G. Savin \cite{Sa88} has given an isomorphism between the
Iwahori-Hecke algebras in the case when $p$ is relatively prime to $2n$ and
$G$ is any reductive $p$-adic group which splits over ${\cal O}$. This
isomorphism generalized a result of Flicker and Kazhdan \cite{FK86} who
considered the special case of $GL(r)$. (Savin's result is discussed briefly
in \S \ref{section:savin} below.) The works of Rallis-Schiffmann \cite{RS77} and
Schultz \cite{Sc98} also dealt with a correspondence between representations on
$\overline{G}$ and on $G$ (or a closely related group), though from a
different perspective and using different methods. Our hope is that the result
in the present paper will form one small step, along with Waldspurger's
multiplicity one theorem for $\overline{SL(2,F)}$ \cite{Wa80} and the
Arthur-Selberg trace formula, in an eventual proof of Labesse's multiplicity
one conjecture for $SL(2)$.

\subsection{Work after 1998}
\label{sec:lit-after}

Since 1998, Hecke algebras of covering groups have been studied with the
theory of types and, more recently, with $L$-groups for covers. This
subsection summarizes the work closest to the present paper.

In the same year as the original paper, Bushnell and Kutzko \cite{BK98}
developed the theory of types. A type is a pair consisting of a compact open
subgroup and a representation of it, whose Hecke algebra controls one
Bernstein component of the category of smooth representations. The algebra
${\cal H}(G//B_k)$ studied here is the Hecke algebra attached to the pair
consisting of $B_k$ and its trivial representation. This paper does not use
the theory of types, and does not address whether that pair is a type.

Savin \cite{Sa04} studied unramified genuine representations of more general
covering groups. Gan and Savin \cite{GS12} used types for the even and odd
Weil representations of the metaplectic group $Mp(2n)$. They proved that the
resulting Hecke algebras are isomorphic to Iwahori--Hecke algebras of odd
orthogonal groups, when the residual characteristic is odd. Takeda and Wood
\cite{TW18} extended this to even residual characteristic. Chen and Li
\cite{CL25} revisited the Gan--Savin correspondence from the point of view of
endoscopy. They showed that it preserves $L$-parameters. These results concern
the Iwahori level, as in Section \ref{section:savin}, rather than the deeper
levels $B_k$ of the present paper.

Karasiewicz \cite{Ka21} gave a Bernstein presentation of a genuine Hecke
algebra on the double cover of a simply laced Chevalley group over
$\mathbb{Q}_2$. He used it to identify the algebra with an Iwahori--Hecke
algebra of a linear group. Peskin \cite{Pe13} studied genuine mod $p$
representations of the metaplectic cover of $SL(2,F)$. She showed that its
genuine mod $p$ spherical Hecke algebra is isomorphic to that of $PGL(2,F)$.

Finally, Weissman \cite{We18} defined an $L$-group for a broad class of
covering groups, the Brylinski--Deligne covers. The volume containing his
paper also has a historical introduction by Gan, Gao and Weissman
\cite{GGW18}. In this framework the dual group of the $n$-fold cover of
$SL(2)$ depends on the parity of $n$, in agreement with the description used
in Section \ref{section:savin}.

\section{Background}
\label{sec:background}

This section explains the concepts the paper relies on, in more detail than
the glossary. Each subsection ends with an example taken from the running
example, Example \ref{ex:running}. Readers who know the theory can skip to
Section \ref{sec:notation}.

\subsection{Hecke algebras of totally disconnected groups}
\label{sec:bg-hecke}

This subsection explains why the generalized Hecke algebra of Definition
\ref{def:1} is the natural object for studying representations. Let $H$ be a
locally compact, totally disconnected, unimodular group, and let $K$ be a
compact open subgroup. A function on $H$ that is locally constant, compactly
supported and bi-$K$-invariant is a finite linear combination of
characteristic functions of double cosets $KhK$. Convolution makes these
functions into an associative algebra, and the characteristic function of $K$
(suitably normalized) is its unit.

The reason for studying this algebra is the following. If $(\pi,V)$ is a
smooth representation of $H$, then the operator
$\pi(f)=\int_H f(h)\pi(h)\,dh$ maps $V$ into the space $V^K$ of $K$-fixed
vectors, for every $f$ in the algebra. So $V^K$ is a module over the algebra.
For a large class of representations, the module $V^K$ determines $V$. The
precise statements are recalled in Proposition \ref{prop:borel}. So a
comparison of Hecke algebras of two groups gives a comparison of their
representations.

\begin{example}
\label{ex:bg-hecke}
Take $H=F^\times$ and $K={\cal O}^\times$. Every element of $F^\times$ is
$\pi^m u$ with $m\in\mathbb{Z}$ and $u\in{\cal O}^\times$, so the double cosets
are the sets $\pi^m{\cal O}^\times$. Let $t_m$ be the characteristic function
of $\pi^m{\cal O}^\times$, and give ${\cal O}^\times$ volume $1$. Then
$t_m*t_{m'}=t_{m+m'}$. So the Hecke algebra is the Laurent polynomial algebra
$\mathbb{C}[t,t^{-1}]$, with $t=t_1$. A smooth character $\chi$ of $F^\times$
that is trivial on ${\cal O}^\times$ acts on its one-dimensional space by
$t\mapsto \chi(\pi)$. This identifies such characters with nonzero complex
numbers.
\end{example}

\subsection{The Hilbert symbol and the Kubota cover}
\label{sec:bg-kubota}

This subsection explains the ingredients of the cover $\overline{G}$ defined
in Section \ref{sec:intro}. The $n^{\text{th}}$ power Hilbert symbol
$(\ ,\ )$ is a map $F^\times\times F^\times\rightarrow\mu_n$. It is
multiplicative in each variable, satisfies $(a,b)(b,a)=1$, and $(a,b)=1$ when
$a$ or $b$ is an $n^{\text{th}}$ power. Kubota \cite{Ku67} used it to write
down the explicit $2$-cocycle $\beta$ of Section \ref{sec:intro}. A
$2$-cocycle is exactly what is needed to make the set $G\times\mu_n$ into a
group with the multiplication rule given there. The group $\overline{G}$ is a
central extension of $G$ by $\mu_n$. It is not a matrix group in general.

\begin{example}
\label{ex:bg-kubota}
In the running example, take $a\in F^\times$ and $h(a)=\operatorname{diag}(a,
a^{-1})$. For diagonal matrices the lower left entry is $0$, so
$x(h(a))=a^{-1}$. For $a,b\in F^\times$ this gives
\[
\beta(h(a),h(b))=\bigl(b^{-1},a^{-1}\bigr)=(b,a).
\]
With $a=-1$ and $b=3$ in $\mathbb{Q}_3$ we get $\beta(h(-1),h(3))=(3,-1)
=(-1,3)^{-1}=-1$. So the cocycle takes the nontrivial value $-1$ on the
diagonal torus. As a second example, $s_1^2=-I$ and $x(s_1)=1$, $x(-I)=-1$.
So $\beta(s_1,s_1)=(-1,-1)$, which equals $1$ in $\mathbb{Q}_3$ because $-1$ is
a unit.
\end{example}

\subsection{The Iwahori subgroup and the Iwahori--Bruhat decomposition}
\label{sec:bg-bruhat}

This subsection explains how the double cosets that index a basis of the Hecke
algebra are organized. The Iwahori subgroup $B$ of Section \ref{sec:notation}
consists of the matrices in $SL(2,{\cal O})$ whose lower left entry lies in
$\pi{\cal O}$. The affine Weyl group $W_a$ is the infinite dihedral group
generated by the images of $s_1$ and $s_2$. Each element $w$ of $W_a$ is a
reduced alternating word in $s_1$ and $s_2$, and its length $\ell(w)$ is the
number of letters. The Iwahori--Bruhat decomposition says that $G$ is the
disjoint union of the double cosets $BwB$, $w\in W_a$. The filtration
subgroups $B_k$ are normal in $B$, so each double coset $B_kgB_k$ is
contained in some $BwB$. Remark \ref{rem:1} shows that
$\operatorname{vol}(B_kbwB_k)=q^{\ell(w)}$ when $\operatorname{vol}(B_k)=1$.

\begin{example}
\label{ex:bg-bruhat}
The elements of $W_a$ of length at most $2$ are $1$, $s_1$, $s_2$, $s_1s_2$
and $s_2s_1$. The element $t_0=s_1^{-1}s_2=h(\pi)$ of Section
\ref{sec:notation} has length $2$. In the running example with $k=1$, the
formula above gives $\operatorname{vol}(B_1s_1B_1)=3$ and
$\operatorname{vol}(B_1s_1s_2B_1)=9$. The index of $B_1$ in $B$ is
$q^2(q-1)=18$. Indeed, $B$ maps onto the upper triangular subgroup of
$SL(2,\mathbb{F}_3)$, which has order $q(q-1)=6$, and the kernel contains
$B_1$ with index $q=3$, coming from the lower left entry modulo $\pi^2$.
\end{example}

\subsection{Genuine representations, induced and produced modules}
\label{sec:bg-reps}

This subsection explains the representation-theoretic notions of Section
\ref{sec:application}. A representation of $\overline{G}$ is
\textbf{genuine} if $\mu_n$ acts through the fixed embedding
$\mu_n\subset\mathbb{C}^\times$. Genuine representations do not factor through
$G$, so they are the new representations that the cover contributes. A smooth
representation is \textbf{admissible} if each space of vectors fixed by a
compact open subgroup is finite dimensional. Its \textbf{matrix coefficients}
are the functions $g\mapsto\langle v^*,\pi(g)v\rangle$.
An irreducible admissible representation is \textbf{supercuspidal} if its
matrix coefficients are compactly supported, and \textbf{square-integrable} if
they are square-integrable. Definitions \ref{def:13} and \ref{def:17} make
these notions precise for $\overline{G}$.

Borel \cite{Bo76} constructed two ways back from Hecke modules to
representations. Given a module $W$ over ${\cal H}(H//K)$, the
\textbf{induced} module is $C_c(H/K)\otimes_{{\cal H}(H//K)}W$. The
\textbf{produced} module is the smooth part of
$\operatorname{Hom}_{{\cal H}(H//K)}(C_c(K\backslash H),W)$. Proposition
\ref{prop:borel} recalls their properties in the genuine setting.

\begin{example}
\label{ex:bg-reps}
Take $n=1$, so that $\overline{G}=G$, and let $\mathbf{1}$ be the trivial
representation of $G$. Its space of $B_k$-fixed vectors is all of
$\mathbb{C}$. A function $f$ in ${\cal H}(G//B_k)$ acts on it by the scalar
$\int_G f(g)\,dg$. So $f_g$ acts by $\operatorname{vol}(B_kgB_k)$. In the
running example with $k=1$, $f_{s_1}$ acts by $3$. The matrix coefficients of
$\mathbf{1}$ are constant functions. They are not compactly supported, and not
square-integrable, since $G$ is not compact. So $\mathbf{1}$ is neither
supercuspidal nor square-integrable.
\end{example}

\section{Some notation}
\label{sec:notation}

This section fixes the matrices, subgroups and Hecke algebras used in the rest
of the paper. It begins with a table of notation. It then defines the
generalized Hecke algebra of $G$ (Definition \ref{def:1}) and its genuine
analog for $\overline{G}$ (Definition \ref{def:2}).

\begin{center}
\renewcommand{\arraystretch}{1.15}
\begin{longtable}{>{$}l<{$} p{0.72\textwidth}}
\textrm{Symbol} & Meaning \\ \hline
F,\ \pi,\ {\cal O},\ q & $p$-adic field, uniformizer, ring of integers,
order of the residue field ${\Bbb F}_q$ \\
\mu_n,\ \theta & $n^{\text{th}}$ roots of unity in $F$, fixed isomorphism
onto $\mu_n({\Bbb C})$ \\
(\ ,\ ) & $n^{\text{th}}$ power Hilbert symbol \\
\beta,\ x(g) & Kubota cocycle and the entry used to define it
(Section \ref{sec:intro}) \\
G,\ \overline{G},\ \rho & $SL(2,F)$, its $n$-fold cover, the projection
$\overline{G}\rightarrow G$ (in Section \ref{sec:application}, $\rho$ also
denotes the right regular representation) \\
\overline{H} & $\rho^{-1}(H)$ for a subgroup $H$ of $G$ \\
s_1,\ s_2,\ t,\ t_0 & the matrices defined below \\
u_\pm(x),\ h(x) & unipotent and diagonal matrices defined below \\
B,\ B_k & Iwahori subgroup and its $k^{\text{th}}$ filtration subgroup \\
{\frak g},\ {\frak b},\ {\frak b}_k & Lie algebras of $G$, $B$, $B_k$ \\
U_k^\pm,\ T_k & unipotent and diagonal parts of $B_k$
(Lemma \ref{lemma:9} and Remark \ref{rem:torus_center}) \\
W_a,\ W & affine and finite Weyl groups (Section \ref{section:sl2}) \\
S & the set $\{s_1,s_2\}\cup B$ indexing the generators \\
{\cal H}(H//K) & generalized Hecke algebra (Definition \ref{def:1}) \\
{\cal H}(\overline{G}//B_k)_{gen} & metaplectic generalized Hecke algebra
(Definition \ref{def:2}) \\
f_g,\ \overline{f}_g & characteristic function of $B_kgB_k$, and its genuine
analog (Section \ref{sec:lemmas}) \\
\beta',\ s,\ s_K & equivalent cocycle, splitting function, Kubota splitting \\
B_k^* & the split copy $\{(b,1)\mid b\in B_k\}$ of $B_k$ in
$\overline{B}_k$ \\
B^{(n)} & the subgroup of $B$ of Remark \ref{rem:torus_center} \\
{\cal H}^\wedge & equivalence classes of irreducible modules of an algebra
${\cal H}$ \\
\end{longtable}
\end{center}

Let
\[
s_1=
\begin{pmatrix}
0 & -1\\
1 & 0
\end{pmatrix},\quad
s_2=
\begin{pmatrix}
0 & -\pi^{-1}\\
\pi & 0
\end{pmatrix},\quad
t=
\begin{pmatrix}
0 & -1\\
\pi & 0
\end{pmatrix},\quad
t_0=
\begin{pmatrix}
\pi & 0 \\
0 & \pi^{-1}
\end{pmatrix}.
\]
Note $t_0=s_1^{-1}s_2$, $t\notin G$, and $ts_1t^{-1}=s_2$. Although
$t\notin G$, conjugation by $t$ defines an automorphism of $G$ of order $2$
(since $t^2=-\pi I$ is central in $GL(2,F)$).

Let
\[
u_+(x)=
\begin{pmatrix}
1 & x\\
0 & 1
\end{pmatrix},\quad
u_-(x)=
\begin{pmatrix}
1 & 0\\
x & 1
\end{pmatrix},\quad
h(x)=
\begin{pmatrix}
x & 0\\
0 & x^{-1}
\end{pmatrix}.
\]

Let
\[
B=\left\{
\begin{pmatrix}
{\cal O} & {\cal O}\\
\pi {\cal O} & {\cal O}
\end{pmatrix}\right\}\cap G
\]
denote the Iwahori subgroup,
\[
{\frak{b}}=\left\{
\begin{pmatrix}
{\cal O} & {\cal O}\\
\pi {\cal O} & {\cal O}
\end{pmatrix}\right\}\cap \frak{g}
\]
denote its Lie algebra, and for $k>0$ let
\[
B_k=\left(I+t^{2k}
\begin{pmatrix}
{\cal O} & {\cal O}\\
\pi {\cal O} & {\cal O}
\end{pmatrix}\right)\cap G
=\left\{
\begin{pmatrix}
a & b\\
\pi c & d
\end{pmatrix}\in G \;\middle|\;
a-1,\, d-1,\, b,\, c \in \pi^k{\cal O}
\right\}
\]
denote the $k^{\text{th}}$ filtration subgroup of $B$. Note that
$t B t^{-1}=B$ and $t B_k t^{-1}=B_k$. Let
\[
{\frak{b}}_k
=\left\{
\begin{pmatrix}
a & b\\
\pi c & -a
\end{pmatrix}
\;\middle|\;
a,b,c\in \pi^{k}{\cal O}\right\}
\]
denote the Lie algebra of $B_k$.

\begin{definition}
\label{def:1}
Let $H$ denote a locally compact unimodular group and let $K$ denote an open
compact subgroup of $H$ with Haar measure on $H$ normalized so that $K$ has
measure one. Let ${\cal H}(H//K)$ denote the space of functions on $H$ which
are locally constant, compactly supported, and bi-$K$-invariant. This forms an
algebra under convolution,
\[
(f*g)(x)=\int_H f(y)g(y^{-1}x)\,dy,
\]
called the \textbf{generalized Hecke algebra}.
\end{definition}

\begin{example}
\label{ex:def1}
Example \ref{ex:bg-hecke} is the case $H=F^\times$, $K={\cal O}^\times$ of
Definition \ref{def:1}. For the case used in this paper, take
$H=G=SL(2,\mathbb{Q}_3)$ and $K=B_1$, as in the running example. The
characteristic functions $f_1$, $f_{s_1}$ and $f_{s_2}$ of $B_1$,
$B_1s_1B_1$ and $B_1s_2B_1$ lie in ${\cal H}(G//B_1)$. The function $f_1$ is
the unit, and $f_{s_1}$ and $f_{s_2}$ are two of the generators of Proposition
\ref{prop:main}.
\end{example}

We shall study the algebra ${\cal H}(G//B_k)$ in Section \ref{section:sl2}.

Recall that a $2$-cocycle $\alpha :H\times H \rightarrow \mu_n$ of a group $H$
is \textbf{trivial} if there is a function $s:H\rightarrow \mu_n$ (called a
\textbf{splitting}) such that $\alpha(h_1,h_2)=s(h_1)s(h_2)s(h_1h_2)^{-1}$,
for $h_1,h_2\in H$. In this case the cover of $H$ defined by $\alpha$ is said
to \textbf{split}, and the map $h \mapsto (h, s(h)^{-1})$ embeds $H$ as a
subgroup of $\overline{H}$. When $\alpha$ is replaced by the equivalent
cocycle $\alpha'(h_1,h_2)=\alpha(h_1,h_2)s(h_1)^{-1}s(h_2)^{-1}s(h_1h_2)$ on
$H\times H$, the splitting becomes trivial on $H$, and $H$ embeds canonically
as $H^*=\{(h,1)\mid h\in H\}\subset \overline{H}$.

\begin{example}
\label{ex:splitting}
For any function $s:H\rightarrow\mu_n$, the cocycle
$\alpha(h_1,h_2)=s(h_1)s(h_2)s(h_1h_2)^{-1}$ is trivial, with splitting $s$.
The cocycle $\alpha'$ obtained from it as above is identically $1$. In the
running example, $1+3\mathbb{Z}_3$ consists of squares in
$\mathbb{Q}_3^\times$, by Hensel's Lemma. So the next paragraph applies with
$k=1$, and shows that $\overline{B}_1$ splits.
\end{example}

By Hensel's Lemma, if $k>2\operatorname{ord}_\pi(n)$ (and in particular for
any $k\ge 1$ when $\gcd(p,n)=1$), we have
$1+\pi^k{\cal O}\subset (F^\times)^n$. For such $k$, every $u\in 1+\pi^k{\cal O}$
is an $n^{\text{th}}$ power in $F^\times$, so the $n^{\text{th}}$ power
Hilbert symbol satisfies $(u,y)=1$ for all $u\in 1+\pi^k{\cal O}$ and all
$y\in F^\times$. By \cite{Ku67}, Theorem 2, the Kubota splitting function
$s_K:G\rightarrow \mu_n$ defined by
\[
s_K\begin{pmatrix} a & b \\ c & d \end{pmatrix}
=
\begin{cases}
(c,d), & \text{if } cd\neq 0 \text{ and } \operatorname{ord}_\pi(c) \text{ is odd},\\
1, & \text{otherwise},
\end{cases}
\]
trivializes $\beta$ on the principal congruence subgroup of level $\pi^k$, and
satisfies $s_K|_{B_k}\equiv 1$ since $d\in 1+\pi^k{\cal O}\subset (F^\times)^n$
for all elements of $B_k$. In fact, the definition of $\beta$ and the Iwahori
factorization $B_k=U_k^- T_k U_k^+$ imply directly that $\overline{B}_k$
splits.

Recall that a complex-valued function $f$ on $\overline{G}$ is called
\textbf{genuine} if $f(g,\zeta)=\theta(\zeta) f(g,1)$, for $g\in G$ and
$\zeta\in \mu_n$. Here $\theta$ is a fixed isomorphism
$\theta:\mu_n\rightarrow \mu_n({\Bbb{C}})$. It being agreed that we may
identify $\mu_n$ and $\mu_n({\Bbb{C}})$, we shall, \emph{from this point on,
drop $\theta$ from the notation.}

\begin{example}
\label{ex:genuine}
In the running example $n=2$ and $\mu_2=\{\pm 1\}$. A function $f$ on
$\overline{G}$ is genuine when $f(g,-1)=-f(g,1)$ for all $g\in G$. So a genuine
function is determined by its values on the elements $(g,1)$, and changes sign
on the other sheet of the cover. A nonzero function that factors through $G$,
such as a nonzero constant, is never genuine when $n\ge 2$.
\end{example}

\begin{definition}
\label{def:2}
Assume that $k>0$ is chosen so large that $1+\pi^k{\cal O}\subset (F^\times)^n$
(so in particular $\overline{B}_k$ splits). Fix a cocycle $\beta'$ equivalent
to $\beta$ such that $\beta'|_{B_k\times B_k}=1$, and identify $B_k$ with the
split subgroup $B_k^*=\{(b,1)\mid b\in B_k\}\subset \overline{B}_k$. Assume
that the Haar measure on $\overline{G}$ is normalized so that $\overline{B}_k$
has measure $1$. Let ${\cal H}(\overline{G}//B_k)_{gen}$ denote the space of
functions on $\overline{G}$ which are genuine, locally constant, compactly
supported, and bi-$B_k^*$-invariant. This is the \textbf{(metaplectic)
generalized Hecke algebra} on $\overline{G}$.
\end{definition}

\begin{example}
\label{ex:def2}
In the running example, $k=1$ satisfies the hypothesis of Definition
\ref{def:2}, because $1+3\mathbb{Z}_3\subset(\mathbb{Q}_3^\times)^2$. The
function $\overline{f}_1$ defined by $\overline{f}_1(b,\zeta)=\zeta$ for
$b\in B_1$, $\zeta=\pm1$, and $\overline{f}_1=0$ off $\overline{B}_1$, is
genuine and bi-$B_1^*$-invariant. It lies in
${\cal H}(\overline{G}//B_1)_{gen}$, and Proposition \ref{prop:main2} shows
that it is the unit of this algebra.
\end{example}

We will study this algebra in Section \ref{section:metaplectic}. Note that it
depends on the splitting of the cocycle on $B_k$. To have a measure on
$\overline{G}$ which fits with the measure on $G$ used in
\S \ref{section:sl2}, we shall also assume that $\mu_n$ has measure $1$.

\section{Some lemmas}
\label{sec:lemmas}

This section collects the facts about double cosets and convolution that the
proofs in Sections \ref{section:sl2} and \ref{section:metaplectic} need. The
main point is Lemma \ref{lemma:genuine}. It shows that the basic product rule
for characteristic functions of double cosets (Lemma \ref{lemma:general})
still holds for their genuine analogs.

Let ${\Bbb{F}}_q={\cal O}/\pi{\cal O}$ denote the residue field of $F$.

\begin{lemma}
\label{lemma:3}
For $k>0$, $B_k/B_{k+1}\cong {\frak{b}}_k/{\frak{b}}_{k+1}\cong {\Bbb F}_q^3$,
as abelian groups. (The second isomorphism is not canonical.)
\end{lemma}

\begin{proof}
The graph of the first isomorphism is
$\{(g\cdot B_{k+1},X+ {\frak{b}}_{k+1}) \mid
(I+X+{\frak{b}}_{k+1})\cap g\cdot B_{k+1}\neq \emptyset\}$
(namely, $I+X \bmod B_{k+1} \mapsto X \bmod {\frak{b}}_{k+1}$ for
$X\in {\frak{b}}_k$). It is also a special case of a general result of Morris
(Proposition 3.4 in \cite{Mo88}). The second isomorphism is due to the fact that
the surjective homomorphism
\[
\begin{pmatrix}
a & b\\
\pi c & -a
\end{pmatrix}\longmapsto
\bigl(\pi^{-k}a \bmod \pi,\; \pi^{-k}b \bmod \pi,\; \pi^{-k}c \bmod \pi\bigr)
\in {\Bbb F}_q^3
\]
has image ${\Bbb F}_q^3$ and kernel ${\frak{b}}_{k+1}$.
\end{proof}

\begin{example}
\label{ex:lemma3}
In the running example with $k=1$, Lemma \ref{lemma:3} says that $B_1/B_2$ is
an elementary abelian group of order $3^3=27$. The element
$u_+(3)=I+X$ of $B_1$, with $X=\left(\begin{smallmatrix}0&3\\0&0
\end{smallmatrix}\right)\in{\frak b}_1$, maps to $(0,1,0)\in{\Bbb F}_3^3$
under the second isomorphism, since $3^{-1}\cdot 3=1$.
\end{example}

The following result is stated as Proposition 2.2 in Chapter 3 of \cite{HM85}.

\begin{lemma}
\label{lemma:general}
In the notation of Definition \ref{def:1}, let $f_h$ denote the characteristic
function of the double coset $K h K$ in $H$. Then $f_1$ is the unit of
${\cal H}(H//K)$. More generally, if
\[
\operatorname{vol}(Kh_1K)\operatorname{vol}(Kh_2K)=\operatorname{vol}(Kh_1h_2K),
\]
then $f_{h_1}*f_{h_2}=f_{h_1h_2}$.
\end{lemma}

\begin{example}
\label{ex:general}
In Example \ref{ex:bg-hecke}, every double coset $\pi^m{\cal O}^\times$ has
volume $1$. So the volume condition of Lemma \ref{lemma:general} always holds,
and the lemma gives $t_m*t_{m'}=t_{m+m'}$. In the running example, Remark
\ref{rem:1} gives $\operatorname{vol}(B_1s_1B_1)\operatorname{vol}(B_1s_2B_1)
=3\cdot3=\operatorname{vol}(B_1s_1s_2B_1)$. So the lemma gives
$f_{s_1}*f_{s_2}=f_{s_1s_2}$. By contrast, the volume condition fails for
$h_1=h_2=s_1$, since $s_1^2=-I\in B$ and $\operatorname{vol}(B_1(-I)B_1)=1\neq
9$. In fact $f_{s_1}*f_{s_1}\neq f_{-I}$. By Lemma \ref{lemma:9}(d) the left
side has integral $9$, while the right side has integral $1$. Relation A(i) of
Proposition \ref{prop:main} gives the correct value.
\end{example}

We need a version of the previous lemma for genuine functions. Assume $\beta'$
is chosen as in Definition \ref{def:2} so that $\beta'|_{B_k\times B_k}=1$.
Since $1+\pi^k{\cal O}\subset (F^\times)^n$, for any $g\in G$ and any
$b_1,b_2\in B_k$ satisfying $b_1 g b_2 = g$, the compatibility of the Kubota
splitting with $1+\pi^k{\cal O}\subset (F^\times)^n$ ensures that
$(b_1,1)(g,1)(b_2,1)=(g,1)$. Thus, for each $g\in G$, there is a unique
genuine bi-$B_k^*$-invariant function $\overline{f}_g$ supported on
$\overline{B}_k(g,1)\overline{B}_k$ satisfying $\overline{f}_g(g,1)=1$;
explicitly, for $b_1,b_2\in B_k$ and $\zeta\in\mu_n$,
\[
\overline{f}_g\bigl((b_1,1)(g,1)(b_2,1)(1,\zeta)\bigr)=\zeta,
\]
and $\overline{f}_g(\overline{x})=0$ for
$\overline{x}\notin \overline{B}_k(g,1)\overline{B}_k$. When $\beta'$ is
additionally normalized on the double coset $B_k g B_k$ so that
$\beta'(b_1,g)\beta'(b_1g,b_2)=1$ for $b_1,b_2\in B_k$, this takes the simple
form
\[
\overline{f}_{g}(x,\zeta)
=
\begin{cases}
\zeta, & x\in B_kgB_k,\\
0, & \text{otherwise}.
\end{cases}
\]
Let
\[
(\phi_1*\phi_2)(g)=
\int_{\overline{G}}\phi_1(x)\phi_2(x^{-1}g)\,dx,
\quad g\in \overline{G}.
\]

Now we prove a genuine analog of Lemma \ref{lemma:general} above.

\begin{lemma}
\label{lemma:genuine}
Assume $\operatorname{vol}(\overline{B}_k)=1$. If $g_1,g_2\in G$ are such that
$\beta'((g_1,1),(g_2,1))=1$ on the double coset representatives and
\[
\operatorname{vol}(B_kg_1B_k)
\operatorname{vol}(B_kg_2B_k)
=\operatorname{vol}(B_kg_1g_2B_k),
\]
then $\overline{f}_{g_1}*\overline{f}_{g_2}=\overline{f}_{g_1g_2}$.
\end{lemma}

\begin{proof}
We compute
\[
\begin{aligned}
(\overline{f}_{g_1}*\overline{f}_{g_2})(g,1)
&=\int_{\overline{G}}
\overline{f}_{g_1}(x)\overline{f}_{g_2}(x^{-1}(g,1))\,dx\\
&=\int_{G}\overline{f}_{g_1}(x,1)\overline{f}_{g_2}(x^{-1}g,1)
\beta'(x,x^{-1}g)^{-1}\,dx.
\end{aligned}
\]
Since $|\overline{f}_{g_i}(x,1)|=f_{g_i}(x)$, the support of this integral is
contained in the support of $f_{g_1}*f_{g_2}$. By Lemma \ref{lemma:general},
$f_{g_1}*f_{g_2}=f_{g_1g_2}$, so
\[
\operatorname{supp}(\overline{f}_{g_1}*\overline{f}_{g_2})
\subset \operatorname{supp}(\overline{f}_{g_1g_2}).
\]
Since $\overline{f}_{g_1}*\overline{f}_{g_2}$ is genuine and bi-$B_k^*$-invariant,
and the space of genuine bi-$B_k^*$-invariant functions supported on a single
double coset $\overline{B}_k(g_1g_2,1)\overline{B}_k$ is one-dimensional
(spanned by $\overline{f}_{g_1g_2}$), we must have
$\overline{f}_{g_1}*\overline{f}_{g_2}=c\overline{f}_{g_1g_2}$ for some
constant $c$. Evaluating at $(g_1g_2,1)$: by the volume condition
$\operatorname{vol}(B_kg_1B_k)\operatorname{vol}(B_kg_2B_k)
=\operatorname{vol}(B_kg_1g_2B_k)$, the map
$B_k\times B_k\times B_k\rightarrow B_kg_1g_2B_k$ given by
$(b_1,b_2,b_3)\mapsto b_1g_1b_2g_2b_3$ has fiber $B_k$ at $g_1g_2$, so for
$x=g_1b_2$ ($b_2\in B_k$) we have $x^{-1}g_1g_2=b_2^{-1}g_2$ and the integrand
at $g=g_1g_2$ is identically $1$ on $g_1B_k$, yielding $c=\operatorname{vol}(g_1B_k)=1$.
\end{proof}

As in the above proof, we have the following expression for the convolution in
terms of the cocycle.

\begin{lemma}
\label{lemma:convolve}
For any genuine functions $\phi_1,\phi_2\in {\cal H}(\overline{G}//B_k)_{gen}$,
we have
\[
(\phi_1*\phi_2)(g,\zeta)
=\zeta\int_{G}\phi_1(x,1)\phi_2(x^{-1}g,1)
\beta'(x,x^{-1}g)^{-1}\,dx,
\]
for any $g\in G$ and $\zeta\in\mu_n$.
\end{lemma}

Note that the generalized metaplectic Hecke algebra depends on the cocycle
$\beta$, so let us denote this dependence temporarily by
${\cal H}_\beta(\overline{G}//B_k)_{gen}$. If the cocycle is changed to an
equivalent one, say to $\beta'$, where
\[
\beta(g,h)=\beta'(g,h)s(g)s(h)s(gh)^{-1}
\]
is the cocycle modified by $s$ and $s:G\rightarrow \mu_n$ is any function
satisfying $s(1)=1$ and $s|_{B_k}\equiv 1$ (for example, the Kubota splitting
\cite{Ku67,KP84}), then the two algebras are isomorphic.

\begin{lemma}
\label{lemma:isomorphic}
Let $\beta,\beta'$ be as above and let $s:G\rightarrow \mu_n$ satisfy
$s(1)=1$. There is a canonical isomorphism
\[
\phi:{\cal H}_{\beta'}(\overline{G}//B_k)_{gen}
\stackrel{\sim}{\longrightarrow}
{\cal H}_{\beta}(\overline{G}//B_k)_{gen},
\]
as algebras, defined by sending $f(g,\zeta)=\zeta f(g,1)$ to
$\phi(f)(g,\zeta)=\zeta s(g) f(g,1)$. In other words, under this mapping, we
have
\[
\phi(f_1*_{\beta'} f_2)(g,\zeta)
=(\phi(f_1)*_{\beta}\phi(f_2))(g,\zeta),
\]
where $*_\beta$ denotes the convolution with respect to the $\beta$ cocycle
\[
(\phi_1*_\beta \phi_2)(g,\zeta)
=\zeta \int_{G}\phi_1(x,1)\phi_2(x^{-1}g,1)
\beta(x,x^{-1}g)^{-1}\,dx
\]
and, similarly, $*_{\beta'}$ denotes the convolution with respect to the
$\beta'$ cocycle.
\end{lemma}

\begin{proof}
Using $\beta(x,x^{-1}g)^{-1}=\beta'(x,x^{-1}g)^{-1}s(x)^{-1}s(x^{-1}g)^{-1}s(g)$,
we have
\[
\begin{aligned}
(\phi(f_1)*_\beta \phi(f_2))(g,\zeta)
&=\zeta \int_G s(x)f_1(x,1)s(x^{-1}g)f_2(x^{-1}g,1)
\beta(x,x^{-1}g)^{-1}\,dx\\
&=\zeta s(g)\int_G f_1(x,1)f_2(x^{-1}g,1)\beta'(x,x^{-1}g)^{-1}\,dx\\
&=\phi(f_1*_{\beta'}f_2)(g,\zeta),
\end{aligned}
\]
as claimed.
\end{proof}

\section{The structure of ${\cal H}(G//B_k)$}
\label{section:sl2}

We want to determine the structure of the generalized Hecke algebra as a
finitely generated algebra with generators and relations. Then we will do the
same for the metaplectic analog and compare the two.

The affine Weyl group $W_a = N_G(T)/(T\cap B)$ of $G$ is an infinite dihedral
group generated by the images of $s_1$ and $s_2$, and the finite Weyl group
$W$ is generated by the image of $s_1$. Note that in $G$, $s_1^2=s_2^2=-I\in B$,
and every element of the group $\langle s_1,s_2\rangle \subset G$ can be
written up to sign $\pm I\in B$ as a reduced alternating word in $s_1$ and
$s_2$ (in particular, $t_0=s_1^{-1}s_2=h(-1)s_1s_2$). We will use the
Iwahori-Bruhat decomposition
\[
G=BW_aB
\]
to determine the structure of the generalized Hecke algebra.

Recall $(\operatorname{Ad} g)(x)=gxg^{-1}$. Let $S=\{s_1,s_2\}\cup B$.

\begin{proposition}
\label{prop:main}
The algebra ${\cal H}(G//B_k)$ is generated by the functions $f_g$, $g\in S$.
Assume that $\operatorname{vol}(B_k)=1$. These elements are subject to the
relations below.

\noindent
\textbf{(A) Relations for the $s_i$'s ($i=1,2$):}
\begin{enumerate}
\item[(i)] $f_{s_1}*f_{s_1}=c_1\sum_x f_x$, where $x$ runs over a set of
representatives of $((-1)\cdot (\operatorname{Ad} s_1)(B_k)\cdot B_k)/B_k$ and
$c_1=\operatorname{vol}(B_k s_1 B_k)=q$,
\item[(ii)] $f_{s_2}*f_{s_2}=c_2\sum_x f_x$, where $x$ runs over a set of
representatives of $((-1)\cdot (\operatorname{Ad} s_2)(B_k)\cdot B_k)/B_k$ and
$c_2=\operatorname{vol}(B_k s_2 B_k)=q$.
\end{enumerate}

\noindent
\textbf{(B) Relations for elements of $B$:}
\begin{enumerate}
\item[(i)] $f_1$ is the identity of ${\cal H}(G//B_k)$,
\item[(ii)] $f_x*f_y=f_{xy}$, for $x,y\in B$,
\item[(iii)] $f_{xg}=f_{gx}=f_g$, for all $x\in B_k,\ g\in G$.
\end{enumerate}

\noindent
\textbf{(C) Mixed relations:}
\begin{enumerate}
\item[(i)] $f_{s_i}*f_x=f_{\operatorname{Ad} s_i(x)}*f_{s_i}$, for
$x\in B\cap (\operatorname{Ad} s_i^{-1})(B)$ ($i=1,2$),
\item[(ii)] $f_{s_1}*f_{u_+(x)}*f_{s_1}=c_1^2 f_{\operatorname{Ad} s_1(u_+(x))}
=c_1^2 f_{u_-(-x)}$, for $x\in {\cal O}^\times$, where
$c_1=\operatorname{vol}(B_k s_1 B_k)$,
\item[(iii)] $f_{s_2}*f_{u_-(-\pi x)}*f_{s_2}=c_2^2 f_{\operatorname{Ad} s_2(u_-(-\pi x))}
=c_2^2 f_{u_+(x)}$, for $x\in {\cal O}^\times$, where
$c_2=\operatorname{vol}(B_k s_2 B_k)$.
\end{enumerate}

\noindent
(The relations C(ii) and C(iii) must be expressed in terms of the generators
$S$, as in \cite{HM85}. Using the Bruhat identity
$u_-(-x)=u_+(-x^{-1})s_1 h(-x) u_+(-x^{-1})$ for $x\in {\cal O}^\times$, the
relation corresponding to C(ii) is
\[
f_{s_1}*f_{u_+(x)}*f_{s_1}
=c_1^2\bigl(f_{u_+(-x^{-1})}*f_{s_1}*f_{h(-x)}*f_{u_+(-x^{-1})}\bigr).
\]
Conjugating by $t$, where $ts_1t^{-1}=s_2$, $tu_+(y)t^{-1}=u_-(-\pi y)$, and
$th(z)t^{-1}=h(z^{-1})$, the relation corresponding to C(iii) is
\[
f_{s_2}*f_{u_-(-\pi x)}*f_{s_2}
=c_2^2\bigl(f_{u_-(\pi x^{-1})}*f_{s_2}*f_{h(-x^{-1})}*f_{u_-(\pi x^{-1})}\bigr),
\]
and similarly $f_{s_2}*f_{u_+(x)}*f_{s_2}=c_2^2 f_{u_-(\pi^2 x)}$ for
$x\in {\cal O}^\times$.)

Furthermore, these relations form a defining set of relations for
${\cal H}(G//B_k)$ as an algebra.
\end{proposition}

\begin{remark}
\label{rem:1}
\begin{enumerate}
\item This proposition has an analog for $G$ replaced by $GL(n)$, proved in
\cite{HM85}, Theorem 2.1, Chapter 3. Unlike the $GL(n)$ case where the length-zero
element $t$ lies in $GL(n,F)$ and normalizes $B_k$, in $SL(2,F)$ the element
$t_0=s_1^{-1}s_2=\operatorname{diag}(\pi,\pi^{-1})$ has length $2$ in $W_a$
and does not normalize $B_k$. However, the outer automorphism
$\operatorname{Ad}(t)$ of $SL(2,F)$ preserves $B$ and $B_k$ and swaps $s_1$
and $s_2$, allowing us to use the Coxeter generators $\{s_1,s_2\}$ directly.
Note that for $t_0=h(-1)s_1s_2$, length additivity gives
$f_{t_0}=f_{h(-1)}*f_{s_1}*f_{s_2}$ and
$f_{t_0^{-1}}=f_{h(-1)}*f_{s_2}*f_{s_1}$, which satisfy
$f_{t_0}*f_{s_1}=f_{s_1}*f_{t_0^{-1}}$.

\item Note that though $S$ is infinite, the set of generators
$\{f_g \mid g\in S\}$ is finite (since $f_b$ for $b\in B$ depends only on the
coset $bB_k\in B/B_k$), so the algebra is indeed finitely generated.

\item Possibly B(iii) is a superfluous relation once restricted to $B/B_k$.

\item Note that $s_1$ and $s_2$ have order $4$ in $G$ ($s_1^2=s_2^2=-I\in B$),
not order $2$ as the analogous permutation matrices in \cite{HM85} do.

\item The volume constants are computed as follows:
\[
\operatorname{vol}(B_k s_1 B_k)
=\operatorname{vol}(B_k s_1 B_k s_1^{-1})/\operatorname{vol}(B_k)
=|U_k^- B_k/B_k|=|U_k^-/(U_k^-\cap B_k)|=|\pi^k{\cal O}/\pi^{k+1}{\cal O}|=q,
\]
and since $tB_kt^{-1}=B_k$ and $ts_1t^{-1}=s_2$,
\[
\operatorname{vol}(B_k s_2 B_k)
=\operatorname{vol}(t B_k s_1 B_k t^{-1})
=\operatorname{vol}(B_k s_1 B_k)=q.
\]
More generally, for any reduced alternating word $w=s_{i_1}s_{i_2}\cdots s_{i_r}$
($i_j\in\{1,2\}$, $i_j\neq i_{j+1}$) and any $b\in B$, Iwahori factorization
for $B_k$ gives
\[
\operatorname{vol}(B_k b w B_k)=\operatorname{vol}(B_k w B_k)
=\prod_{j=1}^r \operatorname{vol}(B_k s_{i_j} B_k)=q^r,
\]
so by Lemma \ref{lemma:general},
$f_w=f_{s_{i_1}}*f_{s_{i_2}}*\cdots *f_{s_{i_r}}$.
\end{enumerate}
\end{remark}

Here are some general facts we will use in the proof below.

\begin{lemma}
\label{lemma:9}
\begin{enumerate}
\item[(a)] $B_k$ is a normal subgroup of $B$, and $tB_kt^{-1}=B_k$.
\item[(b)] Let $U^+_k=\{u_+(x) \mid x\in \pi^k{\cal O}\}$ and
$U^-_k=\{u_-(x) \mid x\in \pi^k{\cal O}\}$. We have
\[
B_k\cdot (s_1^{-1}B_ks_1)=(s_1^{-1}B_ks_1)\cdot B_k=U^-_kB_k,
\]
and
\[
B_k\cdot (s_2^{-1}B_ks_2)=(s_2^{-1}B_ks_2)\cdot B_k=U^+_{k-1}B_k,
\]
and both $U^-_kB_k$ and $U^+_{k-1}B_k$ are groups containing $B_k$ as a
normal subgroup of index $q$.
\item[(c)] In general, we have
\[
(f_{g_1}*f_{g_2})(x)=\int_G f_{g_1}(y)f_{g_2}(y^{-1}x)\,dy
=\operatorname{vol}(B_kg_1B_k\cap xB_kg_2^{-1}B_k).
\]
\item[(d)] In general, we have
\[
\int_G (f*g)(x)\,dx=\int_G f(y)\,dy \int_G g(x)\,dx,\quad
\int_G (f*g*h)(x)\,dx=\int_G f(x)\,dx \int_G g(y)\,dy\int_G h(z)\,dz.
\]
\end{enumerate}
\end{lemma}

\begin{proof}
Parts (a) and (b) may be proven by direct matrix calculations (with the second
identity in (b) following from the first by conjugating by $t$). Parts (c) and
(d) are simple consequences of the definitions.
\end{proof}

\begin{proof}[Proof of Proposition \ref{prop:main}]
We shall follow the ideas in \cite{HM85}, proof of Theorem 2.1.

We begin by verifying A(i): $f_{s_1}*f_{s_1}=c_1\sum_{x\in X} f_x$, where $x$
runs over a complete set $X$ of representatives of
$((-1)\cdot (\operatorname{Ad} s_1)(B_k)\cdot B_k)/B_k$. As in \cite{HM85}, parts
(a) and (b) of Lemma \ref{lemma:9} imply that
\[
f_{s_1}*f_{s_1}=c_1\cdot \chi_{(-1)\cdot U^-_kB_k},
\]
where $\chi_E$ denotes the characteristic function of a subset $E$ of $G$ and
where $c_1\neq 0$ is some constant. In other words,
\[
(f_{s_1}*f_{s_1})(g)=c_1\sum_{x\in X} f_x(g),\quad g\in G.
\]
Plugging $g=-I=s_1^2$ into both sides and using Lemma \ref{lemma:9}(c), we
obtain
\[
\operatorname{vol}(B_k s_1 B_k)=c_1.
\]
Conjugating by the measure-preserving automorphism $\operatorname{Ad}(t)$
(which maps $B_k$ to $B_k$ and $s_1$ to $s_2$) immediately gives A(ii) with
$c_2=\operatorname{vol}(B_k s_2 B_k)=q$.

Part (c) of Lemma \ref{lemma:9} implies part B(i) of the proposition.

We now verify part B(ii). We know $\operatorname{vol}(B_k)=1$ implies
$\operatorname{vol}(bB_k)=1$, for each $b\in B$. Since $B_k$ is normal in $B$,
$B_kbB_k=bB_k$, so Lemma \ref{lemma:general} implies part B(ii).

Part B(iii) follows from the definition of $f_g$.

We shall now verify C(i). Since $B_k$ is normal in $B$, for any $b\in B$ we
have
\[
\operatorname{vol}(B_k s_i b B_k)=\operatorname{vol}(B_k s_i b B_k b^{-1} b)
=\operatorname{vol}(B_k s_i B_k b)=\operatorname{vol}(B_k s_i B_k).
\]
This and Lemma \ref{lemma:general} imply $f_{s_i}*f_x=f_{s_ix}$, for $x\in B$.
Likewise, if $s_ixs_i^{-1}\in B$ then
$f_{s_ixs_i^{-1}}*f_{s_i}=f_{s_ix}$. This implies C(i).

We next verify C(ii). The proof of this part is sketched in \cite{HM85}. The
argument in \cite{HM85} gives $f_{s_1}*f_{u_+(x)}*f_{s_1}=c f_{u_-(-x)}$, for
$x\in {\cal O}^\times$, for some constant $c\neq 0$. Here $c$ may be evaluated
using part (d) of Lemma \ref{lemma:9}, which implies
\[
c=\frac{\operatorname{vol}(B_k s_1 B_k)^2\operatorname{vol}(B_k u_+(x) B_k)}
{\operatorname{vol}(B_k u_-(-x) B_k)}
=\operatorname{vol}(B_k s_1 B_k)^2 = c_1^2.
\]
(Here we have used the fact that $u_+(x)\in B$, so
$\operatorname{vol}(B_k u_+(x) B_k)=\operatorname{vol}(B_k u_+(x))
=\operatorname{vol}(B_k)=1$, and the fact that
$\operatorname{vol}(B_k u_-(-x) B_k)
=\operatorname{vol}(s_1^{-1}B_ks_1 u_+(x) s_1^{-1}B_ks_1)
=\operatorname{vol}(s_1^{-1}B_ks_1 u_+(x))
=\operatorname{vol}(s_1^{-1}B_ks_1)=1$.) Since
$u_-(-x)=u_+(-x^{-1})s_1 h(-x)u_+(-x^{-1})$ with $u_+(-x^{-1}), h(-x)\in B$,
we have $f_{u_-(-x)}=f_{u_+(-x^{-1})}*f_{s_1}*f_{h(-x)}*f_{u_+(-x^{-1})}$ by
C(i) and B(ii).

Part C(iii) follows directly from C(ii) by applying the automorphism
$\operatorname{Ad}(t)$.

It remains to show that this list of relations is a defining set of relations
for the algebra. To this end, we follow the proof of \cite{HM85}, Chapter 3,
Theorem 2.1, as follows.

For $g\in S$, let $\tilde{f}_g$ denote an abstract element. Let
$\tilde{\cal H}$ denote the free algebra generated by these elements modulo
the relations (A)--(C) in the proposition. We want to show that the natural
homomorphism
\begin{equation}
\label{eq:obvious}
\tilde{\cal H}\longrightarrow {\cal H}(G//B_k)
\end{equation}
is an isomorphism.

Choose a set of reduced alternating words $w=s_{i_1}s_{i_2}\cdots s_{i_r}$
($i_j\in\{1,2\}$, $i_j\neq i_{j+1}$) representing the elements of $W_a$, and
define $\tilde{f}_w=\tilde{f}_{s_{i_1}}*\tilde{f}_{s_{i_2}}*\cdots
*\tilde{f}_{s_{i_r}}$ (with $\tilde{f}_1$ for the trivial word). Let
$\tilde{\cal J}$ denote the vector space span in $\tilde{\cal H}$ generated by
the elements $\tilde{f}_g=\tilde{f}_x*\tilde{f}_w*\tilde{f}_y$, where $g=xwy$
for $x,y\in B$ and $w\in W_a$. As in \cite{HM85}, pages 38--39, the relations (B)
and (C) imply that $\tilde{f}_g$ depends only on the double coset $B_kgB_k$,
not on the particular representation $g=xwy$. Furthermore, $\tilde{\cal J}$ is
invariant under left and right multiplication by any abstract generator
$\tilde{f}_g$, $g\in S$: multiplication by $\tilde{f}_b$ ($b\in B$) is absorbed
via B(ii), while multiplication by $\tilde{f}_{s_i}$ either increases the
length of $w$ or is reduced via C(i), C(ii), C(iii), and A(i)--(ii) exactly as
in \cite{HM85}, pages 38--39.

From this invariance property and $\tilde{f}_1\in\tilde{\cal J}$, it follows
that $\tilde{\cal J}=\tilde{\cal H}$. Thus $\tilde{\cal H}$ is spanned by the
elements $\tilde{f}_g$ indexed by the double cosets $B_kgB_k\in B_k\backslash G/B_k$.
Since the images $f_g\in {\cal H}(G//B_k)$ of distinct double cosets have
disjoint supports, they are linearly independent in ${\cal H}(G//B_k)$, which
implies that the map (\ref{eq:obvious}) is injective.

Finally, since $B_k$ is normal in $B$, for any $b_1,b_2\in B$ and $w\in W_a$
we have
\[
\operatorname{vol}(B_k b_1 w b_2 B_k)
=\operatorname{vol}(b_1 B_k w B_k b_2)
=\operatorname{vol}(B_k w B_k).
\]
This and Lemma \ref{lemma:general} imply
$f_{b_1}*f_w*f_{b_2}=f_{b_1wb_2}$. Combining this with the Iwahori-Bruhat
decomposition $G=BW_aB$ and $f_w=f_{s_{i_1}}*\cdots *f_{s_{i_r}}$, we see
that ${\cal H}(G//B_k)$ is spanned by the image of $\tilde{\cal H}$, so the
map (\ref{eq:obvious}) is surjective.
\end{proof}

\begin{example}
\label{ex:prop-main}
In the running example with $k=1$, relation A(i) reads
$f_{s_1}*f_{s_1}=3\sum_x f_x$. By Lemma \ref{lemma:9}(b), the set
$(-1)\cdot U_1^-B_1$ is a union of $q=3$ cosets $xB_1$, so the sum has three
terms. We can check the relation against the trivial representation of Example
\ref{ex:bg-reps}, in which $f_g$ acts by $\operatorname{vol}(B_1gB_1)$. The
left side acts by $3\cdot 3=9$. On the right side each $x$ lies in $B$, so
each $f_x$ acts by $1$, and the right side acts by $3\cdot 3=9$ as well.
\end{example}

\section{The structure of ${\cal H}(\overline{G}//B_k)_{gen}$}
\label{section:metaplectic}

This section is the metaplectic analog of the previous section. We want to
determine the structure of the metaplectic generalized Hecke algebra as a
finitely generated algebra with generators and relations.

Throughout this section, we assume that $k>0$ is chosen so large that
$1+\pi^k{\cal O}\subset (F^\times)^n$ (for instance,
$k>2\operatorname{ord}_\pi(n)$). As noted in Section \ref{sec:notation},
this guarantees that
$(u,y)=1$ for all $u\in 1+\pi^k{\cal O}$ and all $y\in F^\times$, so that
$\overline{B}_k$ splits. If the residual characteristic $p$ of $F$ is
relatively prime to $n$, then $1+\pi{\cal O}\subset (F^\times)^n$ already holds
for $k=1$, and in fact $\overline{B}$ splits.

\begin{remark}
\label{rem:torus_center}
Before stating Proposition \ref{prop:main2}, let us note how the $n$-fold
cover behaves on the diagonal torus $T=\{h(a)\mid a\in F^\times\}$ and on $B$:
\begin{enumerate}
\item For $h(a),h(b)\in T$, Kubota's cocycle gives $\beta(h(a),h(b))=(b,a)$,
so the commutator in $\overline{T}$ is
$[(h(a),1),(h(b),1)]=(b,a)^2=(b^2,a)$.
\item When $n=1$ or $n=2$ (the classical $2$-fold metaplectic cover),
$(b,a)^2=1$ for all $a,b\in F^\times$, so $\overline{T}$ is abelian. Moreover,
when $1+\pi^k{\cal O}\subset (F^\times)^n$, for any $n\ge 1$ the subgroup
$T_k=T\cap B_k=\{h(u)\mid u\in 1+\pi^k{\cal O}\}$ satisfies $(u,a)=1$ for all
$a\in F^\times$, so $\overline{T}_k$ lies in the center of $\overline{T}$.
\item For general $n\ge 1$, let $B^{(n)}=\{b\in B \mid \text{the diagonal entries
of } b \text{ lie in } ({\cal O}^\times)^n(1+\pi^k{\cal O})\}$. On $B^{(n)}$,
all diagonal entries are $n^{\text{th}}$ powers in $F^\times$, so the Hilbert
symbol vanishes identically on $B^{(n)}$ and $\overline{B^{(n)}}$ splits and
commutes with $(t_0,1)$. In addition, choosing coset representatives of
$B/B^{(n)}\cong {\cal O}^\times/({\cal O}^\times)^n$ and of the double cosets
$B_k\backslash G/B_k$, we may choose the coboundary $s:G\rightarrow \mu_n$ in
Lemma \ref{lemma:isomorphic} (with $s|_{B_k}\equiv 1$) so that $\beta'$
trivializes the cocycle on the chosen coset decompositions.
\end{enumerate}
\end{remark}

Let the Haar measure on $\overline{G}$ be normalized so that $\overline{B}_k$
has measure one. Convolution gives ${\cal H}(\overline{G}//B_k)_{gen}$ the
structure of an algebra.

\begin{proposition}
\label{prop:main2}
Assume that $\operatorname{vol}(\overline{B}_k)=1$, that $k>0$ is chosen so
large that $1+\pi^k{\cal O}\subset (F^\times)^n$, and choose the equivalent
cocycle $\beta'$ via Lemma \ref{lemma:isomorphic} as above. The algebra
${\cal H}(\overline{G}//B_k)_{gen}$ is generated by the functions
$\overline{f}_g$, $g\in S$. On the $n^{\text{th}}$-power Iwahori subgroup
$B^{(n)}$ (and on all of $B$ when $n\mid 2$ and $\gcd(p,n)=1$, or under the
normalized coset splitting of Remark \ref{rem:torus_center}), these elements
are subject to the same relations as in Proposition \ref{prop:main}, except
that $f_g$ must be replaced by $\overline{f}_g$. Furthermore, these relations
form a defining set of relations for ${\cal H}(\overline{G}//B_k)_{gen}$ as an
algebra.
\end{proposition}

\begin{proof}
First we verify that the claimed relations hold.

\textbf{A(i), A(ii):} By Lemma \ref{lemma:convolve}, we have
\[
(\overline{f}_{s_i}*\overline{f}_{s_i})(x,1)
=\int_{B_ks_iB_k}
\overline{f}_{s_i}(y^{-1}x,1)\beta'(y,y^{-1}x)^{-1}\,dy.
\]
Because $1+\pi^k{\cal O}\subset (F^\times)^n$, $\beta'$ is trivial on
$U_k^-B_k$ and $U_{k-1}^+B_k$, and for the same reason as in the
non-metaplectic case, $\overline{f}_{s_i}*\overline{f}_{s_i}$ is supported on
$\overline{((-1)\cdot (\operatorname{Ad} s_i)(B_k)\cdot B_k)}$. Evaluating at
$(-I,1)$ gives the same constant $c_i=\operatorname{vol}(B_ks_iB_k)=q$ as in
the non-metaplectic case, which proves A(i) and A(ii).

\textbf{B(i):} For $g\in \overline{G}$ and $\phi\in {\cal H}(\overline{G}//B_k)_{gen}$,
since $\phi(gx^{-1})\overline{f}_1(x)=\phi(g)$ for all $x=(b,\zeta)\in\overline{B}_k$
(by right $B_k^*$-invariance and genuineness of $\phi$), we have
\[
(\phi*\overline{f}_1)(g)=\int_{\overline{G}}\phi(gx^{-1})\overline{f}_1(x)\,dx
=\int_{\overline{B}_k}\phi(gx^{-1})\overline{f}_1(x)\,dx
=\phi(g)\operatorname{vol}(\overline{B}_k)=\phi(g).
\]
Similarly, $\overline{f}_1*\phi=\phi$, so $\overline{f}_1$ is the identity in
${\cal H}(\overline{G}//B_k)_{gen}$.

\textbf{B(ii):} By Lemma \ref{lemma:convolve}, we have
\[
(\overline{f}_x*\overline{f}_y)(g,1)
=\int_{B_kxB_k}
\overline{f}_{y}(h^{-1}g,1)
\beta'(h,h^{-1}g)^{-1}\,dh.
\]
This is non-zero if and only if $gy^{-1}=xb$ for some $b\in B_k$, since $B_k$
is normal in $B$. This implies
$\operatorname{supp}(\overline{f}_x*\overline{f}_y)\subset
\overline{B_kxyB_k}$, and Lemma \ref{lemma:genuine} yields
$\overline{f}_x*\overline{f}_y=\overline{f}_{xy}$.

\textbf{B(iii):} This is an immediate consequence of Lemma \ref{lemma:genuine}.

\textbf{C(i):} The proof of this identity is analogous to the proof in the
non-metaplectic case (using Lemma \ref{lemma:genuine} in place of Lemma
\ref{lemma:general} where appropriate) and is omitted.

\textbf{C(ii), C(iii):} These are proven in the same way as in the
non-metaplectic case, so we only sketch the argument. The proof in the
non-metaplectic case shows that each of these triple convolutions is supported
on a single double $\overline{B}_k$-coset. On the other hand, the support of
each non-zero function in the Hecke algebra contains at least one double
$\overline{B}_k$-coset, since different double $\overline{B}_k$-cosets are
disjoint. Since convolution preserves genuineness and bi-$B_k^*$-invariance,
those functions whose support is precisely one double $\overline{B}_k$-coset
are constant multiples of the corresponding $\overline{f}_g$. The constant
multiple is evaluated at the double coset representative exactly as in the
non-metaplectic case, and since the measures differ only by the normalized
measure on $\mu_n$, the volume computations are identical.

It remains to show that this list of relations is a defining set of relations
for the algebra. Since $\{\overline{f}_g \mid B_kgB_k\in B_k\backslash G/B_k\}$
forms a vector space basis of ${\cal H}(\overline{G}//B_k)_{gen}$ indexed by
the same double cosets as ${\cal H}(G//B_k)$, the argument is identical to the
non-metaplectic case in Proposition \ref{prop:main}.
\end{proof}

\begin{example}
\label{ex:prop-main2}
In the running example, $n=2$ divides $2$ and $\gcd(p,n)=\gcd(3,2)=1$. So
Proposition \ref{prop:main2} applies with $k=1$, and its relations hold on all
of $B$, not only on $B^{(2)}$. For instance, relation A(i) becomes
$\overline{f}_{s_1}*\overline{f}_{s_1}=3\sum_x\overline{f}_x$, with the same
three representatives $x$ as in Example \ref{ex:prop-main}.
\end{example}

\section{Application to representation theory}
\label{sec:application}

This section combines Propositions \ref{prop:main} and \ref{prop:main2} into
the main theorem, an isomorphism of the two Hecke algebras. It then translates
the isomorphism into a correspondence of representations (Corollary
\ref{corollary:main}). Finally it shows that the correspondence preserves
supercuspidal and square-integrable representations (Corollary \ref{cor:18}).
The background for this section is in Section \ref{sec:bg-reps}.

The following is our main result:

\begin{theorem}
\label{thrm:main}
If $k>0$ is chosen sufficiently large as in Proposition \ref{prop:main2} (so
in particular $\overline{B}_k$ splits), then the map
$f_{g}\longmapsto \overline{f}_{g}$, $g\in S$, defines an isomorphism of
algebras ${\cal H}(G//B_k)\rightarrow {\cal H}(\overline{G}//B_k)_{gen}$.
\end{theorem}

\begin{proof}
This is an immediate consequence of Propositions \ref{prop:main} and
\ref{prop:main2}.
\end{proof}

\begin{corollary}
\label{corollary:main}
If $k>0$ is as in Theorem \ref{thrm:main}, then the algebra isomorphism
$\eta:{\cal H}(G//B_k)\rightarrow {\cal H}(\overline{G}//B_k)_{gen}$ induces a
set-theoretic bijection
$\eta^*:{\cal H}(\overline{G}//B_k)_{gen}^\wedge
\rightarrow {\cal H}(G//B_k)^\wedge$ on the sets of equivalence classes of
irreducible (genuine, in the case of $\overline{G}$) representations.
Furthermore, if
\[
\begin{array}{ccc}
{\cal H}(G//B_k)&\longrightarrow &{\cal H}(\overline{G}//B_k)_{gen}\\
f & \longmapsto &\overline{f}
\end{array}
\]
and
\[
\begin{array}{ccc}
{\cal H}(\overline{G}//B_k)_{gen}^\wedge
  &\longrightarrow &{\cal H}(G//B_k)^\wedge\\
\overline{\pi}&\longmapsto &\pi
\end{array}
\]
then $\operatorname{tr}(\overline{\pi}(\overline{f}))=\operatorname{tr}(\pi(f))$.
\end{corollary}

\begin{proof}
Let $\eta:{\cal H}(G//B_k)\rightarrow {\cal H}(\overline{G}//B_k)_{gen}$ be
the algebra isomorphism of Theorem \ref{thrm:main}, and let
$\eta^*:{\cal H}(\overline{G}//B_k)_{gen}^\wedge
\rightarrow {\cal H}(G//B_k)^\wedge$ be the pullback map defined by
$\eta^*(\overline{\pi})=\overline{\pi}\circ \eta$. Suppose
$\pi=\eta^*(\overline{\pi})$ and $\pi'=\eta^*(\overline{\pi}')$ are equivalent
in ${\cal H}(G//B_k)^\wedge$. Since $\eta$ is surjective, every
$\overline{f}\in {\cal H}(\overline{G}//B_k)_{gen}$ is of the form
$\overline{f}=\eta(f)$ for some $f\in {\cal H}(G//B_k)$, so any intertwining
isomorphism between $\pi$ and $\pi'$ also intertwines $\overline{\pi}$ and
$\overline{\pi}'$. Thus $\eta^*$ is injective. Conversely, for any
$(\pi,W)\in {\cal H}(G//B_k)^\wedge$, the composition
$\overline{\pi}=\pi\circ \eta^{-1}$ defines an irreducible representation in
${\cal H}(\overline{G}//B_k)_{gen}^\wedge$ satisfying
$\eta^*(\overline{\pi})=\pi$, so $\eta^*$ is surjective.

The equality of the traces $\operatorname{tr}(\overline{\pi}(\eta(f)))
=\operatorname{tr}(\pi(f))$ holds by definition of $\eta^*$.
\end{proof}

\begin{example}
\label{ex:main}
In the running example with $k=1$, Theorem \ref{thrm:main} gives an
isomorphism ${\cal H}(SL(2,\mathbb{Q}_3)//B_1)\rightarrow
{\cal H}(\overline{G}//B_1)_{gen}$ with $f_g\mapsto\overline{f}_g$. Example
\ref{ex:bg-reps} gives the one-dimensional module of ${\cal H}(G//B_1)$ in
which $f_g$ acts by $\operatorname{vol}(B_1gB_1)$. Under the bijection of
Corollary \ref{corollary:main}, it corresponds to the one-dimensional module
of ${\cal H}(\overline{G}//B_1)_{gen}$ in which each generator
$\overline{f}_g$, $g\in S$, acts by the same number
$\operatorname{vol}(B_1gB_1)$. For example, both $f_{s_1}$ and
$\overline{f}_{s_1}$ act by $3$, so the traces in Corollary
\ref{corollary:main} agree.
\end{example}

\begin{definition}
\label{def:13}
Let $K$ denote a compact open subgroup of $\overline{G}$. A locally constant
complex-valued function $f$ on $\overline{G}$ is called \textbf{left} (resp.,
\textbf{right}) \textbf{Hecke finite} if the vector space spanned by all
functions of the form $h*f$ (resp., $f*h$), for
$h\in {\cal H}(\overline{G}//K)$, is finite dimensional. Let
${\cal A}_{gen}(\overline{G})$ denote the space of all locally constant
complex-valued genuine functions $f$ on $\overline{G}$ for which
\begin{itemize}
\item $f$ is double $K$-finite,
\item $f$ is right- and left-Hecke finite.
\end{itemize}
This space is called the space of \textbf{matrix coefficients} of
$\overline{G}$. If $(\pi,V)$ is a genuine admissible representation of
$\overline{G}$ then the span of the elements
$f(g)=\langle v_0^*,\pi(g)v_0\rangle$ ($g\in \overline{G}$), for all
$v_0\in V$ and $v_0^*\in V^*$, the smooth contragredient of $V$, is called the
space of \textbf{matrix coefficients of $\pi$} and is denoted
${\cal A}_{gen}(\pi)$.
\end{definition}

\begin{example}
\label{ex:def13}
Take $n=1$ and the trivial representation $\mathbf{1}$ of $G$. Its matrix
coefficients $g\mapsto\langle v_0^*,v_0\rangle$ are the constant functions.
So ${\cal A}(\mathbf{1})$ is the one-dimensional space of constants. When
$n\ge 2$ a nonzero constant is not genuine (Example \ref{ex:genuine}), which
reflects the fact that the matrix coefficients of a genuine representation
are genuine functions.
\end{example}

Let $\rho$ (resp., $\lambda$) denote the right (resp., left) regular
representation of $\overline{G}$ on ${\cal A}_{gen}(\overline{G})$. The next
proposition establishes that the metaplectic analogs of (the relevant parts
of) Corollary 1.10.6 in \cite{Si79} for $\overline{G}$ are true. The result
below implies that we can ``explicitly'' realize the space of any genuine
irreducible representation of $\overline{G}$ as a subspace of
${\cal A}_{gen}(\overline{G})$, a fact we will make use of later.

\begin{proposition}
\label{prop:silberger}
Let $f\in {\cal A}_{gen}(\overline{G})$ and let $V_f$ denote the
$\overline{G}$-module spanned by the translates $\rho(g)f$, $g\in \overline{G}$,
and let $\pi$ denote the restriction of $\rho$ to $V_f$. We have:
\begin{enumerate}
\item[(a)] There exist elements $v_0\in V_f$ and $v_0^*\in V_f^*$, the
contragredient of $V_f$, such that $f(g)=\langle v_0^*,\pi(g)v_0\rangle$, for
all $g\in \overline{G}$.
\item[(b)] ${\cal A}_{gen}(\pi)$ is spanned by $\lambda(x)\rho(y)f$,
$x,y\in \overline{G}$.
\end{enumerate}
\end{proposition}

\begin{proof}
The proof of this follows from the general arguments given in \cite{Si79}. The
results in that part of the book \cite{Si79} hold for more general totally
disconnected groups, so one need only check that the argument remains valid
when ``genuineness'' is also assumed.
\end{proof}

We need to recall some facts regarding the relationship between smooth
representations of $\overline{G}$ and finite dimensional
${\cal H}(\overline{G}//B_k)_{gen}$-modules. To this end, we recall briefly
how one can construct smooth representations of $\overline{G}$ from finite
dimensional ${\cal H}(\overline{G}//B_k)_{gen}$-modules---both the so-called
``induced'' and ``produced'' modules in \S 2 of Borel \cite{Bo76} will do this
job. The ``inverse'' construction is then given along with some of the basic
properties of these constructions in Proposition \ref{prop:borel} below.

\begin{definition}
\label{def:15}
If $(r,W)$ is a finite dimensional ${\cal H}(\overline{G}//B_k)$-module then
we define $(I(r),I(W))$ to be the \textbf{induced} $\overline{G}$-module and
$(P(r),P(W))$ to be the (smooth) \textbf{produced} $\overline{G}$-module
constructed in \cite{Bo76}. Analogously, if $(r,W)$ is a finite dimensional
${\cal H}(\overline{G}//B_k)_{gen}$-module then we define
$(I_{gen}(r),I_{gen}(W))$ to be the \textbf{induced} genuine
$\overline{G}$-module and $(P_{gen}(r),P_{gen}(W))$ to be the
\textbf{produced} genuine $\overline{G}$-module obtained by replacing
$C_c(\overline{G}/B_k)$ by $C_c(\overline{G}/B_k)_{gen}$ in the constructions
in \cite{Bo76}.
\end{definition}

\begin{example}
\label{ex:def15}
Take $n=1$, and let $W=\mathbb{C}$ be the module of Example \ref{ex:bg-reps},
on which $f\in{\cal H}(G//B_k)$ acts by $\int_Gf$. It is irreducible, being
one-dimensional. By Proposition \ref{prop:borel}(d) (in the case $n=1$), the
induced module $I(W)$ is irreducible. Take $V=\mathbf{1}$ in Proposition
\ref{prop:borel}(c). Then $V^{B_k}=\mathbb{C}$ is the module $W$ itself, so
$\operatorname{Hom}_G(I(W),\mathbf{1})\cong
\operatorname{Hom}_{{\cal H}(G//B_k)}(W,W)\neq 0$. A nonzero map between
irreducible representations is an isomorphism. So $I(W)$ is isomorphic to the
trivial representation $\mathbf{1}$.
\end{example}

\begin{proposition}
\label{prop:borel}
Let $(r,W)$ be a finite dimensional ${\cal H}(\overline{G}//B_k)_{gen}$-module
and let $(\pi,V)$ be a genuine smooth representation of $\overline{G}$.
\begin{enumerate}
\item[(a)] There are natural ${\cal H}(\overline{G}//B_k)_{gen}$-module
isomorphisms
\[
W\cong I_{gen}(W)^{B_k}\quad \text{and}\quad W\cong P_{gen}(W)^{B_k}.
\]
\item[(b)] $I_{gen}(W)$ is generated, as a $\overline{G}$-module, by
$I_{gen}(W)^{B_k}$.
\item[(c)] The natural restriction maps yield isomorphisms
\[
\rho_P : \operatorname{Hom}_{\overline{G}}(V,P_{gen}(W))
\longrightarrow
\operatorname{Hom}_{{\cal H}(\overline{G}//B_k)_{gen}}(V^{B_k},W)
\]
and
\[
\rho_I : \operatorname{Hom}_{\overline{G}}(I_{gen}(W),V)
\longrightarrow
\operatorname{Hom}_{{\cal H}(\overline{G}//B_k)_{gen}}(W,V^{B_k}).
\]
\item[(d)] $(r,W)$ is an irreducible ${\cal H}(\overline{G}//B_k)_{gen}$-module
if and only if $(I_{gen}(r),I_{gen}(W))$ is an irreducible representation of
$\overline{G}$. Moreover, if $(r,W)$ is an irreducible
${\cal H}(\overline{G}//B_k)_{gen}$-module then $I_{gen}(r)\cong P_{gen}(r)$.
\end{enumerate}
\end{proposition}

\begin{proof}
In the case $n=1$, the statements in the proposition above are either in
\S\S 1--2 of \cite{Bo76} as stated or follow from Borel's results using Schur's
lemma (see, for example, Lemma 3.14 in \cite{CR81}). In general, one need only
check that the argument remains valid when ``genuineness'' is also assumed.
\end{proof}

\begin{definition}
\label{def:17}
We define a genuine admissible irreducible representation of $\overline{G}$ to
be \textbf{supercuspidal} if and only if its matrix coefficients are compactly
supported. We define a genuine admissible irreducible representation of
$\overline{G}$ to be \textbf{square-integrable} if and only if its matrix
coefficients belong to $L^2(\overline{G})$.

We define an irreducible finite dimensional
${\cal H}(\overline{G}//B_k)_{gen}$-module $(r,W)$ to be \textbf{supercuspidal}
(resp., \textbf{square-integrable}) if and only if $I_{gen}(r)$ is a
supercuspidal (resp., square-integrable) representation of $\overline{G}$.
\end{definition}

\begin{example}
\label{ex:def17}
For $n=1$, the trivial representation of $G=SL(2,F)$ is neither supercuspidal
nor square-integrable (Example \ref{ex:bg-reps}). The Steinberg
representation of $SL(2,F)$ is square-integrable but not supercuspidal. It has
a nonzero Iwahori-fixed vector, and a representation with a nonzero
Iwahori-fixed vector is a subquotient of an unramified principal series, so it
is not supercuspidal. By Example \ref{ex:def15}, the trivial
${\cal H}(G//B_k)$-module is neither supercuspidal nor square-integrable in
the sense of the second paragraph of Definition \ref{def:17}.
\end{example}

Next we observe that, as in the case of $GL(n)$, the correspondence of
Corollary \ref{corollary:main} preserves supercuspidals and discrete series
representations (\cite{HM85}, page 33).

\begin{corollary}
\label{cor:18}
Assume $k>0$ is as in Corollary \ref{corollary:main}.
\begin{itemize}
\item An admissible genuine representation
$\overline{\pi}\in {\cal H}(\overline{G}//B_k)_{gen}^\wedge$ is supercuspidal
if and only if $\pi$ is, where $\pi,\overline{\pi}$ are as in Corollary
\ref{corollary:main}.
\item $\overline{\pi}$ is square-integrable if and only if $\pi$ is, where
$\pi,\overline{\pi}$ are as in Corollary \ref{corollary:main}.
\end{itemize}
\end{corollary}

\begin{proof}
Suppose that the mapping
${\cal H}(\overline{G}//B_k)_{gen}^\wedge \rightarrow {\cal H}(G//B_k)^\wedge$
sends a representation $(\overline{r},\overline{W})$ to a representation
$(r,W)$. Note that the construction of the correspondence in Corollary
\ref{corollary:main} implies that $W=\overline{W}$. From the construction of
$I_{gen}(\overline{W})= C_c(\overline{G}/B_k)_{gen}
\otimes_{{\cal H}(\overline{G}//B_k)_{gen}}\overline{W}$ (more precisely, from
the definition of $\otimes_{{\cal H}(\overline{G}//B_k)_{gen}}$), we see that
the $\overline{r}(\overline{f}_g)$-action on
$\overline{W}\cong I_{gen}(\overline{W})^{B_k}$ corresponds to the
$\overline{f}_g$-action (via convolution) on $C_c(\overline{G}/B_k)_{gen}$.
Identifying $\overline{W}$ with $I_{gen}(\overline{W})^{B_k}$, we see that
each element of $\overline{W}$ is Hecke-finite and may therefore be regarded
as a matrix coefficient of $\overline{G}$. Let $(\overline{\pi},\overline{V})$
denote the $\overline{G}$-module generated by an element
$w\in \overline{W}\cong I_{gen}(\overline{W})^{B_k}$. By Proposition
\ref{prop:silberger} (or Proposition \ref{prop:borel}(b)),
$\overline{\pi}=I_{gen}(\overline{r})$.

From this and the fact that $W=\overline{W}$, it follows that if the elements
of $\overline{W}$ are compactly supported (when regarded as matrix
coefficients of $\overline{G}$) then the elements of $W$ are also compactly
supported (when regarded as matrix coefficients of $G$). Now assume that
$(\overline{r},\overline{W})$ is supercuspidal, so the matrix coefficients of
$I_{gen}(\overline{r})$ have compact support. As indicated above, $(r,W)$ is
supercuspidal as well, and the converse holds symmetrically. This establishes
the first part of the corollary.

The second part of the corollary is proved similarly, hence is omitted.
\end{proof}

\begin{example}
\label{ex:cor18}
In the running example, the genuine module of Example \ref{ex:main}
corresponds to the trivial representation of $SL(2,\mathbb{Q}_3)$. By Example
\ref{ex:def17} the trivial representation is neither supercuspidal nor
square-integrable. So Corollary \ref{cor:18} shows that the genuine
representation $I_{gen}(W)$ of $\overline{G}$ attached to this module is
neither supercuspidal nor square-integrable.
\end{example}

\section{The Iwahori algebra}
\label{section:savin}

This section treats the level $k=0$, where $B_k$ is replaced by the Iwahori
subgroup $B$ itself. The method of the earlier sections needs $k>0$. At level
$0$ the genuine Hecke algebra is instead described by a theorem of Savin, which
compares $\overline{G}$ with a linear group that depends on the parity of $n$.

In the case where $k=0$ and $\gcd(p,2n)=1$, the genuine Iwahori-Hecke algebra
${\cal H}(\overline{G}//B)_{gen}$ is described by a general theorem due to
Savin \cite{Sa88}. Recall that for $G=SL(2,F)$, the dual group of the $n$-fold
metaplectic cover $\overline{G}$ depends on the parity of $n$: let
$G^\vee=SL(2,F)$ when $n$ is odd, and let $G^\vee=PGL(2,F)\cong SO(3,F)$ when
$n$ is even (with Iwahori subgroup $B^\vee \subset G^\vee$). We recall Savin's
theorem in this setting:

\begin{theorem}[Savin \cite{Sa88}]
\label{thm:savin}
Assume $\gcd(p,2n)=1$. There is an isomorphism of Iwahori-Hecke algebras
\[
{\cal H}(\overline{G}//B)_{gen}\stackrel{\sim}{\longrightarrow}
{\cal H}(G^\vee//B^\vee),
\]
where $G^\vee=G=SL(2,F)$ when $n$ is odd, and $G^\vee=PGL(2,F)$ when $n$ is
even.
\end{theorem}

\begin{corollary}
\label{cor:20}
Assume $\gcd(p,2n)=1$. The isomorphism
${\cal H}(\overline{G}//B)_{gen}\rightarrow {\cal H}(G^\vee//B^\vee)$ induces
a set-theoretic bijection
${\cal H}(\overline{G}//B)_{gen}^\wedge
\rightarrow {\cal H}(G^\vee//B^\vee)^\wedge$ on the sets of equivalence
classes of irreducible (genuine, in the case of $\overline{G}$)
representations with a $B$-fixed (resp., $B^\vee$-fixed) vector. Furthermore,
if
\[
\begin{array}{ccc}
{\cal H}(G^\vee//B^\vee)&\longrightarrow &{\cal H}(\overline{G}//B)_{gen}\\
f & \longmapsto &\overline{f}
\end{array}
\]
and
\[
\begin{array}{ccc}
{\cal H}(\overline{G}//B)_{gen}^\wedge &\longrightarrow &{\cal H}(G^\vee//B^\vee)^\wedge\\
\overline{\pi}&\longmapsto &\pi
\end{array}
\]
then $\operatorname{tr}(\overline{\pi}(\overline{f}))=\operatorname{tr}(\pi(f))$.
\end{corollary}

\begin{example}
\label{ex:savin}
In the running example, $\gcd(p,2n)=\gcd(3,4)=1$ and $n=2$ is even. So
Theorem \ref{thm:savin} identifies ${\cal H}(\overline{G}//B)_{gen}$ with the
Iwahori--Hecke algebra of $G^\vee=PGL(2,\mathbb{Q}_3)$. In that algebra, with
the Iwahori subgroup of volume $1$, the characteristic function $T_s$ of the
double coset of a simple reflection $s$ satisfies the quadratic relation
$T_s*T_s=(q-1)T_s+qT_1$ of Iwahori and Matsumoto. With $q=3$ this reads
$T_s*T_s=2T_s+3T_1$.
\end{example}

\section{Summary of results}
\label{sec:summary}

This section collects the results of the paper in one place and lists the
questions it leaves open. Let $F$ be a $p$-adic field of characteristic $0$
with $|\mu_n(F)|=n$, let $G=SL(2,F)$, and let $\overline{G}$ be its $n$-fold
Kubota cover.
\begin{enumerate}
\item Proposition \ref{prop:main} gives generators and defining relations for
the generalized Hecke algebra ${\cal H}(G//B_k)$, $k>0$.
\item Proposition \ref{prop:main2} gives the same generators and relations for
the genuine algebra ${\cal H}(\overline{G}//B_k)_{gen}$, under the hypotheses
stated there, when $1+\pi^k{\cal O}\subset(F^\times)^n$.
\item Theorem \ref{thrm:main} deduces that $f_g\mapsto\overline{f}_g$,
$g\in S$, is an algebra isomorphism ${\cal H}(G//B_k)\rightarrow
{\cal H}(\overline{G}//B_k)_{gen}$.
\item Corollary \ref{corollary:main} turns this into a bijection between
irreducible modules of the two algebras that preserves traces. Corollary
\ref{cor:18} shows that the bijection preserves supercuspidal and
square-integrable representations.
\item At the Iwahori level, Theorem \ref{thm:savin} (due to Savin
\cite{Sa88}) and Corollary \ref{cor:20} give the analogous statements, with
$G$ replaced by $G^\vee$, when $\gcd(p,2n)=1$.
\end{enumerate}

Two questions remain open. The first is raised in Remark \ref{rem:1}: whether
relation B(iii) is superfluous once it is restricted to $B/B_k$. The second is
the hope stated in Section \ref{sec:lit-before}: that these results, with
Waldspurger's multiplicity one theorem \cite{Wa80} and the Arthur--Selberg
trace formula, may contribute to a proof of Labesse's multiplicity one
conjecture for $SL(2)$.

\section{Glossary}
\label{sec:glossary}

This section gives short definitions of the technical terms used in the
paper. Section \ref{sec:background} explains most of them in more detail.

\begin{description}
\item[Admissible representation] A smooth representation whose space of
vectors fixed by any compact open subgroup is finite dimensional.
\item[Affine Weyl group] The group $W_a=N_G(T)/(T\cap B)$. For $SL(2,F)$ it is
infinite dihedral, generated by the images of $s_1$ and $s_2$.
\item[Cocycle ($2$-cocycle)] A function $\alpha:H\times H\rightarrow\mu_n$ that
makes $H\times\mu_n$ a group under
$(h_1,\zeta_1)(h_2,\zeta_2)=(h_1h_2,\alpha(h_1,h_2)\zeta_1\zeta_2)$.
\item[Filtration subgroup] The normal subgroup $B_k$ of the Iwahori subgroup
defined in Section \ref{sec:notation}.
\item[Generalized Hecke algebra] The convolution algebra of locally constant,
compactly supported, bi-$K$-invariant functions (Definition \ref{def:1}).
\item[Genuine function or representation] One on which $\mu_n$ acts through
the fixed embedding $\mu_n\subset\mathbb{C}^\times$.
\item[Hilbert symbol] The $n^{\text{th}}$ power norm residue symbol
$F^\times\times F^\times\rightarrow\mu_n$.
\item[Induced and produced modules] Borel's two constructions of a
representation from a Hecke algebra module (Definition \ref{def:15}).
\item[Iwahori subgroup] The subgroup $B$ of $SL(2,{\cal O})$ of matrices whose
lower left entry lies in $\pi{\cal O}$.
\item[Iwahori--Bruhat decomposition] The decomposition $G=BW_aB$ of $G$ into
double cosets of $B$.
\item[Kubota cocycle] The explicit cocycle $\beta$ of Section \ref{sec:intro}
defining the $n$-fold cover of $SL(2,F)$.
\item[Matrix coefficient] A function $g\mapsto\langle v^*,\pi(g)v\rangle$
(Definition \ref{def:13}).
\item[Metaplectic cover] The topological central extension $\overline{G}$ of
$G$ by $\mu_n$ defined by the Kubota cocycle.
\item[Splitting] A function $s$ with
$\alpha(h_1,h_2)=s(h_1)s(h_2)s(h_1h_2)^{-1}$. A cover splits when such an $s$
exists.
\item[Square-integrable representation] One whose matrix coefficients are
square-integrable (Definition \ref{def:17}).
\item[Supercuspidal representation] One whose matrix coefficients are
compactly supported (Definition \ref{def:17}).
\end{description}

\section*{Acknowledgments}

I thank Bill Wardlaw for some helpful conversations on the proof of
Proposition \ref{prop:main} and the anonymous referee for many helpful
corrections and suggestions.

This corrected and expanded version was prepared with the help of three AI
systems. Argon (Google Gemini) acted as reviser. It made the 2026
mathematical corrections to the original text. ChatGPT (OpenAI) acted as
referee of the corrected text. Claude (Anthropic, Claude Opus 5.5) acted as
editor for this draft. It brought the paper into line with the author's style
guide, and drafted the new expository material: the literature review, the
background section, the notation table, the examples, the summary of results
and the glossary. Claude made no changes to the statements or proofs. Review
by AI systems is not an independent certification of the results. Further
details on the AI usage can be found in the revision log,
\texttt{hecke-alg-corresp-for-metapletic-SL2-revision-log.tex}.

\begin{thebibliography}{GGW18}

\bibitem[Bo76]{Bo76} A. Borel, ``Admissible representations of a semi-simple group
over a local field with vectors fixed under an Iwahori subgroup,''
\emph{Invent.\ Math.}\ \textbf{35} (1976), 233--259, MR 56 \#3196.

\bibitem[BK98]{BK98} C. J. Bushnell and P. C. Kutzko, ``Smooth representations of
reductive $p$-adic groups: structure theory via types,'' \emph{Proc.\ London
Math.\ Soc.\ (3)}\ \textbf{77} (1998), 582--634, MR 1643417.

\bibitem[CL25]{CL25} F. Chen and W.-W. Li, ``Spectral transfer for metaplectic
groups.\ II.\ Hecke algebra correspondences,'' preprint,
arXiv:2502.00781, 2025.

\bibitem[CR81]{CR81} C. Curtis and I. Reiner, \emph{Methods of representation
theory, I}, Wiley-Interscience, John Wiley and Sons, New York, 1981,
MR 82i:20001.

\bibitem[Fl80]{Fl80} Y. Flicker, ``Automorphic forms on covering groups of
$GL(2)$,'' \emph{Invent.\ Math.}\ \textbf{57} (1980), 119--182, MR 81m:10057.

\bibitem[FK86]{FK86} Y. Flicker and D. Kazhdan, ``Metaplectic correspondence,''
\emph{Publ.\ Math.\ IHES}\ \textbf{64} (1986), 53--110, MR 88d:11049.

\bibitem[GGW18]{GGW18} W. T. Gan, F. Gao and M. H. Weissman, ``$L$-groups and
the Langlands program for covering groups: a historical introduction,''
\emph{Ast\'erisque} \textbf{398} (2018), 1--31.

\bibitem[GS12]{GS12} W. T. Gan and G. Savin, ``Representations of metaplectic
groups II: Hecke algebra correspondences,'' \emph{Represent.\ Theory}\
\textbf{16} (2012), 513--539.

\bibitem[Ge76]{Ge76} S. Gelbart, \emph{Weil's representation and the spectrum of the
metaplectic group}, Lecture Notes in Math., No.\ 530, Springer-Verlag, Berlin,
1976, MR 54 \#12654.

\bibitem[HM85]{HM85} R. Howe (with the collaboration of A. Moy),
\emph{Harish-Chandra homomorphisms for $p$-adic groups}, CBMS Regional
Conference Series in Mathematics, No.\ 59, Amer.\ Math.\ Soc., Providence,
1985, MR 87h:22023.

\bibitem[Ka21]{Ka21} E. Karasiewicz, ``A Hecke algebra on the double cover of
a Chevalley group over $\mathbb{Q}_2$,'' \emph{Algebra Number Theory}\
\textbf{15} (2021), 1729--1753.

\bibitem[KP84]{KP84} D. Kazhdan and S. Patterson, ``Metaplectic forms,''
\emph{Publ.\ Math.\ IHES}\ \textbf{59} (1984), 35--142, MR 85g:22033.

\bibitem[Ku67]{Ku67} T. Kubota, ``Topological covering of $SL(2)$ over a local
field,'' \emph{J.\ Math.\ Soc.\ Japan}\ \textbf{19} (1967), 114--121,
MR 34 \#4264.

\bibitem[Mo88]{Mo88} L. Morris, ``$P$-cuspidal representations,'' \emph{Proc.\
London Math.\ Soc.}\ \textbf{57} (1988), 329--356, MR 89j:22038.

\bibitem[Pe13]{Pe13} L. Peskin, ``Ordinary mod $p$ representations of the
metaplectic cover of $p$-adic $SL_2$,'' PhD thesis, California Institute of
Technology, 2013.

\bibitem[RS77]{RS77} S. Rallis and G. Schiffmann, ``Repr\'esentations
supercuspidales du groupe m\'etaplectique,'' \emph{J.\ Math.\ Kyoto Univ.}\
\textbf{17} (1977), 567--603, MR 58 \#16523.

\bibitem[Sa88]{Sa88} G. Savin, ``Local Shimura correspondence,'' \emph{Math.\ Ann.}\
\textbf{280} (1988), 185--190, MR 89h:22018.

\bibitem[Sa04]{Sa04} G. Savin, ``On unramified representations of covering
groups,'' \emph{J.\ Reine Angew.\ Math.}\ \textbf{566} (2004), 111--134,
MR 2039325.

\bibitem[Sc98]{Sc98} J. Schultz, ``Lifting of characters of $\overline{SL_2(F)}$
and $SO_{2,1}(F)$, for $F$ nonarchimedean,'' PhD thesis (supervised by
J. Adams), Univ.\ of Maryland at College Park, 1998.

\bibitem[Si79]{Si79} A. Silberger, \emph{Introduction to harmonic analysis on
reductive $p$-adic groups}, Princeton Univ.\ Press, Princeton, 1979,
MR 81m:22025.

\bibitem[TW18]{TW18} S. Takeda and A. Wood, ``Hecke algebra correspondences
for the metaplectic group,'' \emph{Trans.\ Amer.\ Math.\ Soc.}\ \textbf{370}
(2018), 1101--1121.

\bibitem[Wa80]{Wa80} J.-L. Waldspurger, ``Correspondance de Shimura,'' \emph{J.\
Math.\ Pures Appl.}\ \textbf{59} (1980), 1--132, MR 83f:10029.

\bibitem[We18]{We18} M. H. Weissman, ``$L$-groups and parameters for covering
groups,'' \emph{Ast\'erisque} \textbf{398} (2018), 33--186.

\end{thebibliography}

\end{document}
